\ssr{ВВЕДЕНИЕ}

Задача коммивояжёра является одной из классических задач комбинаторной оптимизации.
Её актуальность обусловлена широким применением в логистике, маршрутизации и планировании перемещений, особенно в условиях, когда стоимость перемещения между пунктами несимметрична.

Целью данной лабораторной работы является исследование и сравнение двух подходов к решению задачи коммивояжёра на ориентированном графе: точного метода полного перебора и приближённого метаэвристического метода --- муравьиного алгоритма.

В рамках работы поставлены следующие задачи:
\begin{enumerate}
	\item описать метод полного перебора, указать его преимущество и недостатки, представить схему и реализацию алгоритма;
	\item описать муравьиный алгоритм без использования элитных муравьёв, указать его преимущество и недостатки, представить его схему и реализацию;
	\item выполнить оценку трудоёмкости обоих алгоритмов на основе их схем;
	\item провести параметризацию муравьиного алгоритма;
	\item сформировать рекомендации по выбору значений параметров для заданного класса данных, состоящего как минимум из трёх полносвязных ориентированных графов;
	\item провести сравнительный анализ рассмотренных методов с точки зрения качества решений, устойчивости и вычислительной сложности;
\end{enumerate}

Решение задачи выполняется в постановке поиска гамильтонова цикла минимальной длины на ориентированном графе, моделирующем реальные условия воздушной навигации.